\section{Анализ требований к разрабатываемой системе}

На основе анализа задач физической симуляции (раздел~1.1) и обзора существующих решений (раздел~1.2) сформулированы требования к физическому движку «Модейлера».
Требования разделены на функциональные (что система должна делать) и нефункциональные (как она должна это делать).

\subsection{Функциональные требования}

\paragraph{Моделирование твёрдых тел.}
Система должна поддерживать симуляцию произвольного числа твёрдых тел с заданными массово-инерционными характеристиками.
Каждое тело имеет шесть степеней свободы (три поступательные и три вращательные).
Поддерживается применение внешних сил и моментов.

\paragraph{Обнаружение столкновений.}
Система должна обнаруживать геометрические пересечения между телами и вычислять контактные точки, нормали и глубину проникновения.
Поддерживаемые примитивные геометрии: сфера, капсула, коробка (AABB/OBB), цилиндр, выпуклый многогранник.

\paragraph{Физика контактов и трения.}
Обработка контактов включает: нормальные импульсы (предотвращение взаимопроникновения) и касательные импульсы (кулоново трение).
Коэффициенты трения и упругости задаются на уровне пар материалов.

\paragraph{Кинематические связи (joints).}
Поддерживаются следующие типы сочленений: вращательное (revolute), призматическое (prismatic), фиксированное (fixed), шаровое (ball-and-socket).
Ограничения скорости, позиции и момента задаются на уровне сочленения.

\paragraph{Интеграция с ROS2.}
Движок должен предоставлять интерфейс для передачи данных о состоянии тел (поза, скорость, усилия) в систему ROS2 и приёма управляющих воздействий.

\paragraph{GPU-ускорение.}
Система должна поддерживать перенос вычислительно-ёмких этапов симуляции на GPU через API Vulkan Compute.
GPU-бэкенд должен быть опциональным: при его отсутствии система работает в CPU-режиме.

\paragraph{Параллельная симуляция.}
Движок должен поддерживать одновременную симуляцию нескольких независимых сцен (для задач обучения с подкреплением).
На GPU это реализуется через пакетную (batched) обработку сцен.

\subsection{Нефункциональные требования}

\paragraph{Производительность.}
Система должна обеспечивать симуляцию сцены с $N_b \leq 1000$ твёрдых тел в реальном времени ($\Delta t = 1\text{ мс}$) на типичной рабочей станции (8-ядерный CPU).
При использовании GPU время шага для пакета из 1000 независимых сцен с $N_b = 64$ телами должно быть не более 10~мс.

\paragraph{Детерминированность.}
При одинаковых начальных условиях и последовательности управляющих воздействий результаты симуляции должны быть воспроизводимы.
Для CPU-режима детерминированность гарантируется при фиксированном числе потоков; для GPU-режима — в пределах одного запуска на одном устройстве.

\paragraph{Точность.}
Физическая ошибка (отклонение от аналитического решения) для базовых тестовых сценариев (свободное падение, маятник, упругий удар) не должна превышать 1\% за 10 секунд симуляции при шаге интегрирования 1~мс.

\paragraph{Платформенная независимость.}
Движок реализуется на C++17 и должен компилироваться без модификаций под Linux x86-64 и Windows x86-64.
GPU-модуль использует Vulkan~1.3 как кроссплатформенный API, поддерживаемый на оборудовании AMD, NVIDIA и Intel.

\paragraph{Расширяемость.}
Архитектура должна допускать замену отдельных компонентов (солвера, алгоритма обнаружения столкновений) без изменения публичного API.
Это позволяет экспериментировать с альтернативными алгоритмами и адаптировать движок под конкретные задачи.

\paragraph{Интеграция в «Модейлер».}
Физический движок предоставляется как статически компонуемая библиотека (lib) с заголовочными файлами C++.
Зависимости ограничены стандартной библиотекой C++, математической библиотекой линейной алгебры и Vulkan SDK.

\subsection{Выводы по разделу}

Сформулированные требования охватывают все аспекты физического движка: базовую симуляционную функциональность, GPU-ускорение, точность и интеграцию с платформой.
Требования к производительности указывают на необходимость эффективной GPU-реализации и пакетной обработки сцен.
Требование кроссплатформенности и расширяемости определяет выбор Vulkan в качестве GPU API вместо CUDA, что отличает разрабатываемый движок от большинства аналогов.
